\chapter{Introducción}
En \verb|\thesisfaculty{}| se escribe el nombre de la facultad a la que pertenece el
autor.

Para tesis de grado, en \verb|\thesistype{}| se escribe ``Carrera de'' y la carrera a la
que pertenece el autor, por ejemplo \verb|\thesistype{Carrera de Computación}|. Para
tesis de maestría, en \verb|\thesistype{}| se escribe ``Maestría en'' y el nombre de la
maestría, por ejemplo \verb|\thesistype{Maestría en Ciencias de la Computación}|.

En \verb|\thesistitle{}| se escribe el título de la tesis el cual no debe llevar punto
final. En \verb|\thesisdegree{}| se escribe el nombre del título a obtener.

En \verb|\thesisauthor{}| se escriben primero los nombres y luego los apellidos del
autor, sin comas ni punto final. En caso de contar con segundo autor utilizar
\verb|\thesissecondauthor{}|, y si no, comentar o borrar el macro.

En \verb|\thesisdirector{}{}| y \verb|\thesisseconddirector{}{}| se escriben los nombres
de los directores seguidos de su ORCID, se cumplen las mismas normas de los autores, no
se debe agregar el título académico (Ing., Dr., Lic., etc.).

En \verb|\thesisyear{}|, \verb|\thesismonth{}| y \verb|\thesisday{}| se escribe la fecha
de entrega expresada únicamente con números.

En \verb|\thesisresumen{}| se escribe el resumen en español, y en
\verb|\thesisabstract{}| se el mismo en inglés. Ambos deben ser un único párrafo con una
extensión máxima de 250 palabras. En \verb|\resumenkeywords{}| y
\verb|\abstractkeywords{}| se escriben entre tres y cinco palabras clave, separadas por
comas, sin punto final y en minúscula.

Se recomienda organizar el contenido en archivos separados dentro de \verb|content/|, y
mantener una jerarquía máxima de tres niveles: \verb|\chapter{}|, \verb|\section{}| y
\verb|\subsection{}|.

Los agradecimientos son opcionales, en general para secciones antes del contenido
principal se recomienda usar el prefijo \verb|0_| para una mejor organización a nivel de
archivos, por ejemplo \verb|0_1_agradecimientos.tex|. Para el contenido principal se
recomienda usar el prefijo \verb|1_|, por ejemplo \verb|1_1_introduccion.tex|.

Las referencias bibliográficas se manejan con el archivo \verb|references.bib|. Para
citar una fuente se usa \verb|\cite{}|, por ejemplo
\verb|\cite{a, b}| produciría \cite{gao_object_2025, sourav_cnn_2023}.

Los anexos se escriben en el archivo destinado para ello y se crean con el comando
\verb|\anexo{}|. Cada anexo puede llevar su propio título y etiqueta para citarse dentro
del documento, por ejemplo Anexo~\ref{anexo:ejemplo1}.

Se recomienda usar \verb|tabularx| para las tablas (ejemplo la
Tabla~\ref{tab:ejemplo-tabularx}) para ajustar automáticamente el ancho de las columnas
al ancho total de la página, se puede utilizar L, C o R para definir el alineamiento de
cada columna.

\begin{table}[H]
  \centering
  \caption{Ejemplo de tabla con \texttt{tabularx}.}
  \label{tab:ejemplo-tabularx}
  \begin{tabularx}{\textwidth}{LLL}
    \hline
    \textbf{Encabezado1} & \textbf{Encabezado2} & \textbf{Encabezado3} \\
    \hline
    Valor1               & Valor2               & Valor3               \\
    Valor4               & Valor5               & Valor6               \\
    \hline
  \end{tabularx}
\end{table}

Para figuras (ejemplo la Figura~\ref{fig:comparacion-desempeno}), utilizar
\verb|\figurecaption{}{}| para escribir un pie de figura corto seguido de una
descripción larga, unicamente el corto aparecerá en la tabla de figuras.

\begin{figure}[H]
  \centering
  \includegraphics[width=0.2\textwidth]{example-image-a}
  \figurecaption{%
    Comparación de desempeño por época entre un modelo con \texttt{freeze=None} y otro
    con \texttt{freeze=10}.%
  }{%
    La adaptación completa de la red en promedio otorga un mejor desempeño (0.87) a lo
    largo del entrenamiento, frente a mantener las 10
    primeras capas fijas (0.83).%
  }
  \label{fig:comparacion-desempeno}
\end{figure}
